\section{Introduction}

Lung adenocarcinoma (LUAD) is the most common subtype of non-small cell lung cancer and
exhibits substantial molecular heterogeneity.
Transcriptomic profiling of tumour tissue by the Cancer Genome Atlas (TCGA) identified three
reproducible molecular subtypes: the terminal respiratory unit (TRU), characterised by
differentiation markers and enrichment for \textit{EGFR} mutations; the proximal-inflammatory
(PI) subtype, dominated by immune-related gene expression; and the proximal-proliferative (PP)
subtype, marked by cell cycle and \textit{KRAS}-associated programmes \cite{tcga2014luad}.
These subtypes describe the tumour \emph{state} --- the molecular profile of established
disease.
Matched tumour-normal cohorts, such as the \textit{Clinical Proteomic Tumour Analysis Consortium
LUAD dataset} \cite{gillette2020proteogenomic}, were designed to support a complementary
analysis: by profiling adjacent normal lung tissue from the same patient, they make it
possible to attribute molecular differences to the tumour itself rather than to inter-individual
variation in normal tissue composition.
In practice, matched designs are most often used for conventional paired differential expression,
which identifies genes that shift on average across the cohort but discards patient-level
information about how far and in which direction any individual tumour has moved.

The central question of this project is whether patients with LUAD undergo qualitatively
distinct molecular transformations during tumourigenesis, or whether transformation is
essentially uniform across the cohort.
We address this by operating in a multi-omics latent space rather than in the original
high-dimensional feature space.
Multi-Omics Factor Analysis (MOFA+) \cite{argelaguet2020mofa2} is fit jointly on all
tumour and normal samples, yielding a set of latent factors that capture shared axes of
variation across transcriptomic, proteomic, and phosphoproteomic data.
Every sample receives a factor score vector~$Z$; for each patient~$i$ we define the
\emph{latent transformation vector}
\begin{equation}
    \Delta Z_i \;=\; Z_{\text{tumour},i} \;-\; Z_{\text{normal},i},
    \label{eq:deltaz}
\end{equation}
which encodes the direction and magnitude of that patient's tumourigenesis as a displacement
in the shared latent space.
Clustering patients by~$\Delta Z$ thus groups those who undergo qualitatively similar
multi-omics reprogramming, regardless of their absolute tumour state.
We refer to these groups as \emph{transformation subtypes} to distinguish them from
the state-based subtypes identified by prior transcriptomic classification.

This report describes the full MOSAIC pipeline.
We first fit a MOFA+ model on all 204 samples to learn a compact, biologically grounded
latent representation (\S\ref{sec:methods}).
We then compute per-patient~$\Delta Z$ vectors and characterise their geometry --- magnitude,
factor composition, and relationships to clinical covariates.
We cluster the~$\Delta Z$ vectors to identify candidate transformation subtypes and assess
their robustness across alternative representations.
We interpret each subtype biologically using pathway analysis on both MOFA factor loadings
and per-patient differential expression, and finally test preliminary clinical associations
through survival and driver mutation analyses.
